% GENERATED FILE. Edit canonical lessons, then run npm run generate:books.
% canonical-source-hash: fnv1a64:74decf61c33b7ed7
% canonical-lessons: SA-C172-yatri, SA-C172-yanapatram, SA-C172-bhramanam, SA-C172-paryatanam, SA-C172-paryatakah, SA-R172-first-pass-words-from-chapters-168-169, SA-R172-second-pass-words-from-chapters-170-172

\chapter{Passenger, Ticket, Tour, Tourism, Tourist}
\label{ch:sa-passenger-ticket-tour-tourism-tourist}
\hlchaptermodality{172}

\begin{chapteropening}
\textbf{By the end of this chapter:} \emph{I can say a passenger, a ticket, a tour, tourism and a tourist.}

Complete the last lesson of chapter 172: I can say a passenger, a ticket, a tour, tourism and a tourist.
\end{chapteropening}

\section[\texorpdfstring{\sk{यात्री} (yātrī)}{यात्री (yātrī)}]{\sk{यात्री} (yātrī) — a passenger}

\label{lesson:SA-C172-yatri}



\begin{quote}
Before the new one: say the Sanskrit for a programme, then the Sanskrit for a competition.
\end{quote}

\subsection*{You'll want to know: \sk{यात्री}}
\textbf{\sk{यात्री}} — \emph{yātrī} — \textquotedblleft{}a passenger\textquotedblright{}.

\begin{grammarlens}[title={What you've built}]
One more word for this chapter.
\end{grammarlens}

\subsection*{Your turn}
\begin{itemize}
\raggedright
  \item \emph{Say it:} \emph{yātrī}
  \item \emph{Say it:} \emph{yātrī}, once more
  \item \emph{Say it:} say \emph{yātrī}
  \item \emph{Recall:} say the Sanskrit for a programme, then the Sanskrit for a competition, then say \emph{yātrī} again
\end{itemize}

\begin{tcolorbox}[breakable,colback=teal!4,colframe=teal!35!black,title={Before you move on}]
What does \sk{यात्री} mean? (\textquotedblleft{}A passenger\textquotedblright{}.) Say it once more.
\end{tcolorbox}
\section[\texorpdfstring{\sk{यानपत्रम्} (yānapatram)}{यानपत्रम् (yānapatram)}]{\sk{यानपत्रम्} (yānapatram) — a ticket}

\label{lesson:SA-C172-yanapatram}



\begin{quote}
Before the new one: say the Sanskrit for a competition, then the Sanskrit for a passenger.
\end{quote}

\subsection*{You'll want to know: \sk{यानपत्रम्}}
\textbf{\sk{यानपत्रम्}} — \emph{yānapatram} — \textquotedblleft{}a ticket\textquotedblright{}.

\begin{grammarlens}[title={What you've built}]
One more word for this chapter.
\end{grammarlens}

\subsection*{Your turn}
\begin{itemize}
\raggedright
  \item \emph{Say it:} \emph{yānapatram}
  \item \emph{Say it:} \emph{yānapatram}, once more
  \item \emph{Say it:} say \emph{yānapatram}
  \item \emph{Recall:} say the Sanskrit for a competition, then the Sanskrit for a passenger, then say \emph{yānapatram} again
\end{itemize}

\begin{tcolorbox}[breakable,colback=teal!4,colframe=teal!35!black,title={Before you move on}]
What does \sk{यानपत्रम्} mean? (\textquotedblleft{}A ticket\textquotedblright{}.) Say it once more.
\end{tcolorbox}
\section[\texorpdfstring{\sk{भ्रमणम्} (bhramaṇam)}{भ्रमणम् (bhramaṇam)}]{\sk{भ्रमणम्} (bhramaṇam) — a tour, travel}

\label{lesson:SA-C172-bhramanam}



\begin{quote}
Before the new one: say the Sanskrit for a passenger, then the Sanskrit for a ticket.
\end{quote}

\subsection*{You'll want to know: \sk{भ्रमणम्}}
\textbf{\sk{भ्रमणम्}} — \emph{bhramaṇam} — \textquotedblleft{}a tour, travel\textquotedblright{}.

\begin{grammarlens}[title={What you've built}]
One more word for this chapter.
\end{grammarlens}

\subsection*{Your turn}
\begin{itemize}
\raggedright
  \item \emph{Say it:} \emph{bhramaṇam}
  \item \emph{Say it:} \emph{bhramaṇam}, once more
  \item \emph{Say it:} say \emph{bhramaṇam}
  \item \emph{Recall:} say the Sanskrit for a passenger, then the Sanskrit for a ticket, then say \emph{bhramaṇam} again
\end{itemize}

\begin{tcolorbox}[breakable,colback=teal!4,colframe=teal!35!black,title={Before you move on}]
What does \sk{भ्रमणम्} mean? (\textquotedblleft{}A tour, travel\textquotedblright{}.) Say it once more.
\end{tcolorbox}
\section[\texorpdfstring{\sk{पर्यटनम्} (paryaṭanam)}{पर्यटनम् (paryaṭanam)}]{\sk{पर्यटनम्} (paryaṭanam) — tourism}

\label{lesson:SA-C172-paryatanam}



\begin{quote}
Before the new one: say the Sanskrit for a ticket, then the Sanskrit for a tour.
\end{quote}

\subsection*{You'll want to know: \sk{पर्यटनम्}}
\textbf{\sk{पर्यटनम्}} — \emph{paryaṭanam} — \textquotedblleft{}tourism\textquotedblright{}.

\begin{grammarlens}[title={What you've built}]
One more word for this chapter.
\end{grammarlens}

\subsection*{Your turn}
\begin{itemize}
\raggedright
  \item \emph{Say it:} \emph{paryaṭanam}
  \item \emph{Say it:} \emph{paryaṭanam}, once more
  \item \emph{Say it:} say \emph{paryaṭanam}
  \item \emph{Recall:} say the Sanskrit for a ticket, then the Sanskrit for a tour, then say \emph{paryaṭanam} again
\end{itemize}

\begin{tcolorbox}[breakable,colback=teal!4,colframe=teal!35!black,title={Before you move on}]
What does \sk{पर्यटनम्} mean? (\textquotedblleft{}Tourism\textquotedblright{}.) Say it once more.
\end{tcolorbox}
\section[\texorpdfstring{\sk{पर्यटकः} (paryaṭakaḥ)}{पर्यटकः (paryaṭakaḥ)}]{\sk{पर्यटकः} (paryaṭakaḥ) — a tourist}

\label{lesson:SA-C172-paryatakah}



\begin{quote}
Before the new one: say the Sanskrit for a tour, then the Sanskrit for tourism.
\end{quote}

\subsection*{You'll want to know: \sk{पर्यटकः}}
\textbf{\sk{पर्यटकः}} — \emph{paryaṭakaḥ} — \textquotedblleft{}a tourist\textquotedblright{}.

\begin{grammarlens}[title={What you've built}]
One more word for this chapter.
\end{grammarlens}

\subsection*{Your turn}
\begin{itemize}
\raggedright
  \item \emph{Say it:} \emph{paryaṭakaḥ}
  \item \emph{Say it:} \emph{paryaṭakaḥ}, once more
  \item \emph{Say it:} say \emph{paryaṭakaḥ}
  \item \emph{Recall:} say the Sanskrit for a tour, then the Sanskrit for tourism, then say \emph{paryaṭakaḥ} again
\end{itemize}

\begin{tcolorbox}[breakable,colback=teal!4,colframe=teal!35!black,title={Before you move on}]
What does \sk{पर्यटकः} mean? (\textquotedblleft{}A tourist\textquotedblright{}.) Say it once more.
\end{tcolorbox}
\section[(dialogue)]{Words from chapters 168-169 — a first pass}

\label{lesson:SA-R172-first-pass-words-from-chapters-168-169}



\begin{quote}
Say the words for tourism and a tourist.
\end{quote}

\subsection*{Your turn}
Say these aloud:

\begin{itemize}
\raggedright
  \item a class, a chapter, a syllable, a colour, a word
  \item a result, a mark, a row, a dictionary, writing
\end{itemize}

\begin{tcolorbox}[breakable,colback=teal!4,colframe=teal!35!black,title={Before you move on}]
A class? (\textbf{\sk{कक्षा}}, \emph{kakṣā}.) A chapter? (\textbf{\sk{अध्यायः}}, \emph{adhyāyaḥ}.) A syllable? (\textbf{\sk{अक्षरम्}}, \emph{akṣaram}.) A colour? (\textbf{\sk{वर्णः}}, \emph{varṇaḥ}.) A word? (\textbf{\sk{पदम्}}, \emph{padam}.) A result? (\textbf{\sk{परिणामः}}, \emph{pariṇāmaḥ}.) A mark? (\textbf{\sk{अंकः}}, \emph{aṅkaḥ}.) A row? (\textbf{\sk{श्रेणी}}, \emph{śreṇī}.) A dictionary? (\textbf{\sk{शब्दकोशः}}, \emph{śabdakośaḥ}.) Writing? (\textbf{\sk{लेखनम्}}, \emph{lekhanam}.)
\end{tcolorbox}
\section[(dialogue)]{Words from chapters 170-172 — a second pass}

\label{lesson:SA-R172-second-pass-words-from-chapters-170-172}



\begin{quote}
Say the words for work and an action.
\end{quote}

\subsection*{Your turn}
Say these aloud:

\begin{itemize}
\raggedright
  \item work, an action, an occupation, a livelihood, service
  \item a message, a convenience, an arrangement, a programme, a competition
  \item a passenger, a ticket, a tour, tourism, a tourist
\end{itemize}

\begin{tcolorbox}[breakable,colback=teal!4,colframe=teal!35!black,title={Before you move on}]
Work? (\textbf{\sk{कार्यम्}}, \emph{kāryam}.) An action? (\textbf{\sk{कर्म}}, \emph{karma}.) An occupation? (\textbf{\sk{उद्योगः}}, \emph{udyogaḥ}.) A livelihood? (\textbf{\sk{वृत्तिः}}, \emph{vṛttiḥ}.) Service? (\textbf{\sk{सेवा}}, \emph{sevā}.) A message? (\textbf{\sk{सन्देशः}}, \emph{sandeśaḥ}.) A convenience? (\textbf{\sk{सुविधा}}, \emph{suvidhā}.) An arrangement? (\textbf{\sk{व्यवस्था}}, \emph{vyavasthā}.) A programme? (\textbf{\sk{कार्यक्रमः}}, \emph{kāryakramaḥ}.) A competition? (\textbf{\sk{प्रतियोगिता}}, \emph{pratiyogitā}.) A passenger? (\textbf{\sk{यात्री}}, \emph{yātrī}.) A ticket? (\textbf{\sk{यानपत्रम्}}, \emph{yānapatram}.) A tour? (\textbf{\sk{भ्रमणम्}}, \emph{bhramaṇam}.) Tourism? (\textbf{\sk{पर्यटनम्}}, \emph{paryaṭanam}.) A tourist? (\textbf{\sk{पर्यटकः}}, \emph{paryaṭakaḥ}.)
\end{tcolorbox}
